\documentclass[11pt,a4paper]{article}
\usepackage[utf8]{inputenc}
\usepackage[vietnamese]{babel}
\usepackage{amsmath,amssymb}
\usepackage{geometry}
\geometry{a4paper, total={180mm,260mm}, left=15mm, right=15mm, top=18mm, bottom=18mm}

\usepackage{tikz}
\usetikzlibrary{shapes.geometric, arrows.meta, positioning, calc}
\tikzset{
  arr/.style={-{Stealth[scale=0.8]}},
  every node/.append style={align=center}
}
\usepackage{tcolorbox}
\tcbuselibrary{skins,breakable}
\usepackage{fontawesome}
\usepackage{fancyhdr}
\usepackage{enumitem}

\pagestyle{fancy}
\fancyhf{}
\renewcommand{\headrulewidth}{0.5pt}
\renewcommand{\footrulewidth}{0.5pt}
\fancyhead[L]{\small\bfseries\color{violet!85!black}\faBook\ SỔ TAY KIẾN THỨC TOÁN 2}
\fancyhead[R]{\small\bfseries\color{violet!85!black}Lớp Toán Thầy Lĩnh -- loptoanthaylinh.vn}
\fancyfoot[C]{\small\color{gray}-- Sổ tay kiến thức: Bài 30. Ngày - tháng --}

\setlength{\parindent}{0pt}
\setlength{\parskip}{0pt}

\begin{document}

% --- BANNER BÀI HỌC ---
\begin{tcolorbox}[
  enhanced,
  colback=violet!8!white,
  colframe=violet!70!black,
  arc=3mm, boxrule=1.2pt,
  left=3mm, right=3mm, top=2mm, bottom=2mm,
  halign=center
]
  {\LARGE\bfseries\color{violet!90!black}\faBookmark\ BÀI 30: NGÀY -- THÁNG}\\[1mm]
  {\itshape\footnotesize\color{black!75} Các tháng trong năm, số ngày của từng tháng và mẹo đốt ngón tay đếm ngày trên 1 trang A4}
\end{tcolorbox}

\vspace{-2mm}

\noindent
\begin{minipage}[t]{0.49\textwidth}\vspace{0pt}
\begin{tcolorbox}[
  enhanced,
  colback=violet!4!white,
  colframe=violet!75!black,
  arc=2mm, boxrule=0.8pt,
  left=2mm, right=2mm, top=1.8mm, bottom=1.8mm,
  title={\footnotesize\bfseries\color{white}\faKey\ 1. Các tháng trong năm}
]
{\footnotesize
\begin{itemize}\setlength{\itemsep}{1pt}\setlength{\leftskip}{-2.5mm}
  \item Một năm có \textbf{12 tháng}, từ tháng $1$ đến tháng $12$.
  \item \textbf{Các tháng có 31 ngày (7 tháng):}\\
  Tháng $1$, Tháng $3$, Tháng $5$, Tháng $7$, Tháng $8$, Tháng $10$, Tháng $12$.
  \item \textbf{Các tháng có 30 ngày (4 tháng):}\\
  Tháng $4$, Tháng $6$, Tháng $9$, Tháng $11$.
  \item \textbf{Tháng đặc biệt:}\\
  Tháng $2$ có \textbf{28 ngày} (năm thường) hoặc \textbf{29 ngày} (năm nhuận).
\end{itemize}
}
\end{tcolorbox}
\end{minipage}
\hfill
\begin{minipage}[t]{0.49\textwidth}\vspace{0pt}
\begin{tcolorbox}[
  enhanced,
  colback=teal!4!white,
  colframe=teal!75!black,
  arc=2mm, boxrule=0.8pt,
  left=2mm, right=2mm, top=1.8mm, bottom=1.8mm,
  title={\footnotesize\bfseries\color{white}\faSliders\ 2. Kĩ năng xem tờ lịch tháng}
]
{\footnotesize
\begin{itemize}\setlength{\itemsep}{1pt}\setlength{\leftskip}{-2.5mm}
  \item \textbf{Cấu trúc tờ lịch tháng:}\\
  - Hàng trên ghi các \textbf{Thứ} trong tuần: Thứ Hai, Thứ Ba, Thứ Tư, Thứ Năm, Thứ Sáu, Thứ Bảy, Chủ nhật.\\
  - Các ô bên dưới ghi các \textbf{Ngày} trong tháng bằng chữ số.
  \item \textbf{Quy tắc gióng hàng:}\\
  Nhìn ngày cần tìm, gióng thẳng lên hàng trên cùng để biết ngày đó rơi vào \textbf{thứ mấy}.
  \item Mỗi tuần lễ có đúng \textbf{7 ngày}.
\end{itemize}
}
\end{tcolorbox}
\end{minipage}

\vspace{-1.5mm}
\begin{tcolorbox}[
  enhanced,
  colback=orange!4!white,
  colframe=orange!80!black,
  arc=2mm, boxrule=0.8pt,
  left=2mm, right=2mm, top=1.5mm, bottom=1.5mm,
  title={\footnotesize\bfseries\color{white}\faSitemap\ 3. Mẹo ghi nhớ số ngày trong tháng bằng đốt mu bàn tay}
]
\centering
\begin{tikzpicture}[x=1.3cm, y=0.8cm]
  \node[rectangle, rounded corners=2pt, draw=red!80!black, fill=red!15, text width=5cm, align=center, font=\scriptsize\bfseries] at (0,0) {
    \textbf{KHỚP NHÔ LÊN TRÊN (31 NGÀY)}\\Tháng 1, Tháng 3, Tháng 5, Tháng 7,\\Tháng 8, Tháng 10, Tháng 12
  };
  \node[rectangle, rounded corners=2pt, draw=blue!80!black, fill=blue!15, text width=5cm, align=center, font=\scriptsize\bfseries] at (5.5,0) {
    \textbf{KHE LÕM XUỐNG DƯỚI}\\$\bullet$ Tháng 2: $28$ hoặc $29$ ngày.\\$\bullet$ Tháng 4, 6, 9, 11: Có \textbf{30 ngày}.
  };
\end{tikzpicture}
\end{tcolorbox}

\vspace{-1.5mm}
\noindent
\begin{minipage}[t]{0.49\textwidth}\vspace{0pt}
\begin{tcolorbox}[
  enhanced,
  colback=green!4!white,
  colframe=green!60!black,
  arc=2mm, boxrule=0.8pt,
  left=2mm, right=2mm, top=1.8mm, bottom=1.8mm,
  title={\footnotesize\bfseries\color{white}\faCheckCircle\ 4. Ví dụ thực hành chuẩn mực}
]
{\footnotesize
\textbf{1. Trả lời câu hỏi ngày tháng:}\\
- Tháng $8$ có bao nhiêu ngày? $\implies$ Có \textbf{31 ngày}.\\
- Tháng $9$ có bao nhiêu ngày? $\implies$ Có \textbf{30 ngày}.
\textbf{2. Tính thứ trong tuần:}\\
Nếu ngày $5$ tháng $10$ là Thứ Ba, thì ngày $12$ tháng $10$ là thứ mấy?\\
\textit{Giải thích:} $12 - 5 = 7$ ngày (đúng $1$ tuần lễ) $\implies$ Ngày $12$ tháng $10$ cũng là \textbf{Thứ Ba}.
}
\end{tcolorbox}
\end{minipage}
\hfill
\begin{minipage}[t]{0.49\textwidth}\vspace{0pt}
\begin{tcolorbox}[
  enhanced,
  colback=red!4!white,
  colframe=red!75!black,
  arc=2mm, boxrule=0.8pt,
  left=2mm, right=2mm, top=1.8mm, bottom=1.8mm,
  title={\footnotesize\bfseries\color{white}\faExclamationTriangle\ 5. Mẹo tránh bẫy trắc nghiệm}
]
{\footnotesize
\begin{itemize}\setlength{\itemsep}{1pt}\setlength{\leftskip}{-2.5mm}
  \item \textbf{Cặp tháng liền nhau cùng có 31 ngày:}\\
  Tháng \textbf{7} và Tháng \textbf{8} là hai tháng liền kề duy nhất trong mùa hè cùng có \textbf{31 ngày}.
  \item \textbf{Cộng thêm 7 ngày cùng thứ:}\\
  Muốn biết ngày cùng thứ ở tuần sau, ta chỉ việc lấy ngày hiện tại cộng thêm $7$.
\end{itemize}
}
\end{tcolorbox}
\end{minipage}


\end{document}
